\subsection*{Reading an action described in other words}
We held the Chinese reader and image writer fixed and trained an English image encoder to recognize descriptions of an action. With dictionary definitions in the questions, the new system found both object positions and wrote both registered answers on 44 of 268 paired pages, compared with 11 for its parent and six when definition supervision was exchanged among training examples (Fig.~\ref{fig:visual-action-meaning}). On the 241 pairs in which both definitions contained several words, the corresponding counts were 42, seven and three. On the earlier 188 pairs with alternate action names, the counts were 32, seven and eight. Literal-name performance on the original 1,691 pages was 716, 703 and 710. Many English words in the held phrases had individually received optical supervision.

The model completed one of nine pairs of Tang lines, writing \textit{cup} and \textit{child}. Both lines of this pair occurred in the inherited training material. Twelve of the 18 verse instances matched a complete training line, and no pair passed the requirement that both lines and both target windows be unseen. We therefore evaluated the frozen models on ten additional pairs from five stories by Pu Songling. All ten passed the overlap screen. For pictured definitions, the new reader selected both correct lines throughout on four pairs, compared with one for the parent and zero for exchanged supervision. All three wrote both registered answers on zero pairs. When the correct object positions were supplied, the shared writer completed two pairs. For a request about fetching a peach, the model chose the correct line and wrote \textit{you}; with the position of the peach supplied, it wrote \textit{peach}. These ten pairs locate errors in both source selection and lexical writing. They provide limited evidence about the frequency of either error across literary Chinese.
